\section{A transfer lemma for a Gauss--Seidel tail}
\label{sec:revisit-relative-tail}
This section gives a reusable handoff criterion, not a claim that the old
AESP burn-in satisfies it. The relevant mass deficit is relative to an RPPR
optimum. It need not require capturing almost all of the full PPR mass.

\begin{proposition}[Relative-deficit thresholded coordinate tail]
\label{prop:revisit-relative-tail}
Let $0<\rho<\epsppr$ and suppose a prefix returns
$0\leq\overline\vx\leq\vu:=\vx_\rho^*$.
Write $\Delta=\vomega^\trans(\vu-\overline\vx)$ and initialize the actual
PPR residual $\bm R=\vb-\mQ\overline\vx$ by one scan of the sparse prefix
support and its incident entries. While some
$R_i>\alpha\epsppr\omega_i$, make the coordinate update
\begin{equation}
 x_i\gets x_i+\frac{R_i-\alpha\rho\omega_i}{Q_{ii}}.
 \label{eq:revisit-reserved-gs}
\end{equation}
Update that row's residual effects and repeat. Then the tail terminates,
returns $0\leq\vx\leq\vu$, and satisfies
$\|\mD^{-1/2}(\vx-\vx_0^*)\|_\infty\leq\epsppr$.
Its repeated adjacency work is at most
\begin{equation}
 W_{\rm tail}\leq
 \frac{(1+\alpha)\Delta}{2\alpha(\epsppr-\rho)}.
 \label{eq:revisit-relative-tail-work}
\end{equation}
Initialization costs $O(1+1/\rho)$ incidence and scalar operations before
label-map logarithms. If the incoming PPR residual is nonnegative, that
property is preserved and the output has the ACL residual certificate.
\end{proposition}
\begin{proof}
Maintain $\vx\leq\vu$. If $u_i=0$, Stieltjes signs imply
$R_i\leq(\vb-\mQ\vu)_i\leq\alpha\rho\omega_i$, so $i$ is never
selected. If $u_i>0$, complementarity and the off-diagonal signs give
\[
 R_i-\alpha\rho\omega_i
   =[\mQ(\vu-\vx)]_i\leq Q_{ii}(u_i-x_i).
\]
The positive update therefore preserves the order and increases the
weighted primal mass by more than
$\alpha(\epsppr-\rho)d_i/Q_{ii}$.
The total available increase is $\Delta$, and
$Q_{ii}=(1+\alpha)/2$. This proves the work bound and finite termination.
Every selected row belongs to $\supp\vu$, and the updated residual at its
own coordinate is exactly $\alpha\rho\omega_i$; neighbor residuals
increase. This also proves preservation of nonnegative residuals when
that property holds initially.

At termination $\bm R\leq\alpha\epsppr\vomega$ on represented labels.
Outside the seed and the boundary of the actual support, it is zero.
Inverse positivity yields
$\vx_0^*-\vx=\mQ^{-1}\bm R\leq\epsppr\vomega$;
the left side is nonnegative by the maintained order.
Initializing the residual opens only $\supp\overline\vx$, whose volume is
at most $1/\rho$. A FIFO active queue with deterministic label maps charges
each new activation and update to its processed incidence; boundary rows
are opened only if selected. Output materialization costs another
$O(1+1/\rho)$ words.
\end{proof}

The reserve $\alpha\rho\omega_i$ is part of the update. It is equivalent
to a residual-dependent under-relaxation factor
$1-\alpha\rho\omega_i/R_i$. At $\rho=\epsppr/2$, every selected factor
lies in $(1/2,1)$. The proof is not for a fixed over-relaxed SOR sweep.

\begin{proposition}[A directly observable residual-excess budget]
\label{prop:revisit-excess-tail}
More generally, suppose only $0\leq\overline\vx\leq\vx_0^*$ and define
\[
 M_\rho(\vx)=\sum_i\omega_i
       [(\vb-\mQ\vx)_i-\alpha\rho\omega_i]_+.
\]
The same reserved coordinate tail preserves $0\leq\vx\leq\vx_0^*$,
terminates with semantic error at most $\epsppr$, and has adjacency work
\begin{equation}
 W_{\rm tail}\leq
 \frac{Q_{ii}M_\rho(\overline\vx)}
      {\alpha^2(\epsppr-\rho)}.
 \label{eq:revisit-excess-tail}
\end{equation}
Initialization and output cost the volume of the supplied prefix support
and its boundary; that volume must be charged to the prefix separately.
Nonnegative input residuals remain nonnegative.
For $\rho=\epsppr/2$, the local certificate
$M_\rho(\overline\vx)\leq\alpha\theta$ gives at most
$2/(\epsppr\sqrt\alpha)$ tail incidences.
If the stronger order $\overline\vx\leq\vx_\rho^*$ holds, then also
\[
 \Delta\leq M_\rho(\overline\vx)/\alpha.
\]
\end{proposition}
\begin{proof}
For a selected coordinate let
$a=(R_i-\alpha\rho\omega_i)/Q_{ii}>0$ be its increment.
Its own contribution to $M_\rho$ falls by $\omega_iQ_{ii}a$.
The increase of the positive part at a neighbor is at most the increase
of its residual, so all neighbor contributions increase by at most
$ca\omega_i$, where $c=(1-\alpha)/2$.
Therefore
\[
 M_\rho(\vx+a\eunit i)
 \leq M_\rho(\vx)-\alpha\omega_i a.
\]
Since $a>\alpha(\epsppr-\rho)\omega_i/Q_{ii}$, summing this
nonnegative potential proves \cref{eq:revisit-excess-tail} and finite
termination. The PPR order is preserved because
$R_i=[Q(\vx_0^*-\vx)]_i\leq Q_{ii}((x_0^*)_i-x_i)$.
At termination the same inverse-positive residual argument as above gives
the semantic bound. Residual nonnegativity is preserved by the same row
update. The positive residual excess is supported on the source, the
prefix support and its observed boundary, so its sum is locally computable.

For the final assertion set $\ve=\vx_\rho^*-\overline\vx\geq0$
and $\vh=[\vb-Q\overline\vx-\alpha\rho\vomega]_+$.
Complementarity on the optimal support, and the Stieltjes signs off it,
give $Q\ve\leq\vh$. Hence
$\ve\leq Q^{-1}\vh$ and
$\Delta\leq\vomega^\trans\vh/\alpha=M_\rho/\alpha$.
\end{proof}

This potential subtracts the RPPR reserve before summing positive residual
mass. It is different from requiring the full remaining PPR mass to be
small. It also permits a prefix above the RPPR optimum on some coordinates,
provided it remains below PPR. The new ramp itself does not maintain that
PPR order, as the counterexamples show; an arbitrary ramp iterate cannot
be fed into this proposition without a separate certificate.

\begin{corollary}[An objective certificate supplies the relative deficit]
\label{cor:revisit-relative-gap}
If the prefix additionally has certified RPPR gap at most $G$, then
\[
 \Delta\leq\sqrt{\frac{2G}{\alpha\rho}}.
\]
With $\rho=\epsppr/2$ and $G\leq\alpha\mu\rho/2$, the tail costs at
most $2/(\epsppr\sqrt\alpha)$ adjacency entries, apart from its
initialization and bookkeeping charges.
\end{corollary}
\begin{proof}
The difference is supported on $\supp\vu$, of volume at most $1/\rho$.
Weighted Cauchy--Schwarz and strong convexity give
$\Delta\leq\rho^{-1/2}\|\vu-\overline\vx\|_2
\leq\sqrt{2G/(\alpha\rho)}$.
Thus $\Delta\leq\theta\leq\sqrt\alpha$, and
\cref{eq:revisit-relative-tail-work} proves the claim.
\end{proof}

\paragraph{Rounded tail.}
For a dyadic prefix, refine its common density grid to a dyadic $h_{\rm t}$
with $h_{\rm t}\leq\alpha(\epsppr-\rho)/(2Q_{ii})$.
Round only the positive density increment in
\cref{eq:revisit-reserved-gs} downward to this grid. It retains at least
half of the ideal increment, preserves $\vx\leq\vu$, and leaves the
updated density residual below the activation threshold.
The work bound merely doubles. The residual-excess potential also decreases
by at least $\alpha\omega_i$ times the actual rounded increment, so
\cref{eq:revisit-excess-tail} likewise only doubles.
Keep exact degree-weighted residual
numerators: their common denominator uses the fixed input denominator of
$\alpha$ and the density grid. Each unweighted threshold comparison uses
only the individual degree. No product of encountered degree denominators
accumulates. This supplies a bounded-arithmetic tail for a bounded-arithmetic
prefix, with the same local work scale.

\paragraph{Scope of the transfer.}
The criterion applies to any deterministic or certified randomized prefix
that supplies both lower order to $\vx_\rho^*$ and the stated objective
gap, with that prefix's own work and success guarantees. It does not prove
those properties for an unchanged old AESP trajectory.
It also does not, by itself, give a practical speedup when composed with
our generic clipping-based OP2 routine: at this small objective budget,
that routine's clipping already satisfies
\cref{cor:revisit-positive-residual}, so the tail need not do any work.
The lemma is useful for investigating a cheaper, differently certified
prefix and for separating relative RPPR deficit from full-PPR mass capture.
